\section{A literal tangent-cone coefficient forces finite strict HS progress}

This chapter concerns genuine point blowup charts. It does not infer an
initial-diagram inclusion, transformed standardness, or the full all-vertex
order criterion from exceptional saturation. The obstruction is a literal
coefficient of a prescribed original canonical vertex polynomial.

\begin{theorem}[Pure coefficients of the same canonical normalized quotient]
\label{hir:actual-point-canonical-quotient-pure-coefficient}
\uses{hir:actual-original-canonical-point-initial-polynomial-congruence}
\lean{Hironaka.HironakaActualPointTangentCoefficientHSDrop.canonical_point_quotient_pure_coefficient}
\leanok
For an original analytic germ ideal $J$, let $G_a$ be its actual canonical
reduced vertex generator, of degree $d_a$. In a genuine point chart with
pivot $j$ and direction $b$, put
$p_a=\delta_{j,b}(\operatorname{in}_{d_a}G_a)$, using the literal complex
polynomial dehomogenization. Every actual quotient
$e^{d_a}H_a=\pi^*G_a$ satisfies
$[X^\gamma]\widehat H_a=[X^\gamma]p_a$ for every exponent with
$\gamma_j=0$.
\end{theorem}
\begin{proof}
The genuine same-generator exceptional initial congruence gives
$H_a-\operatorname{ev}(p_a)\in(e)$. Its whole Taylor image is divisible
by $X_j$, because the original Taylor algebra map sends this same
coordinate germ to $X_j$. Every coefficient with pivot exponent zero
therefore vanishes in that difference. The original Taylor image of
the actual polynomial evaluation is literally the included polynomial,
by the algebra-map coordinate equations. Subtracting coefficients
proves the claimed identity for the same quotient.
\end{proof}

\begin{theorem}[An internally measured source minimum gives a strict finite drop]
\label{hir:actual-point-low-tangent-coefficient-strict-hs-drop}
\uses{hir:actual-point-canonical-quotient-pure-coefficient,hir:actual-original-point-normalized-exceptional-regularity,hir:actual-hs-profile-point-power-detection,hir:actual-original-canonical-reduced-vertex-basis}
\lean{Hironaka.HironakaActualPointTangentCoefficientHSDrop.hilbertSamuelValue_point_normalized_lt_of_low_dehom_coefficient}
\leanok
Let $\mu$ be the infimum in $\mathbb N$ of the degrees of the actual
canonical initial vertices of $J$. Suppose some original vertex $a$
and pure exponent $\gamma_j=0$ have
$r=|\gamma|<\mu$ and $[X^\gamma]p_a\ne0$.
Then the entire independently defined point-normalized ideal $N$ has
$H_N(r)<H_J(r)$. No normalized quotient, order certificate or diagram
comparison is an input. The empty vertex family has $\mu=0$ and hence
cannot satisfy the disclosed coefficient witness.
\end{theorem}
\begin{proof}
Each canonical original generator has no Taylor coefficients below its
own vertex degree. All vertex degrees are at least $\mu>r$, so each
generator belongs to the actual maximal-ideal power $\mathfrak m^{r+1}$.
Their genuine whole span is $J$, giving $J\subseteq\mathfrak m^{r+1}$.
Thus the actual finite-jet HS value $H_J(r)$ equals the ambient value.

Internally choose the genuine canonical normalized family once. Its
member $H_a$ lies in the independently defined whole normalized ideal.
The preceding theorem and the nonzero literal coefficient show
$H_a\notin\mathfrak m^{r+1}$. Consequently that whole ideal is not
contained in this power. The genuine finite-jet kernel/rank criterion
makes its HS value different from the ambient value, while the natural
surjection from the ambient finite jet gives the upper bound by that
value. These two facts give strict inequality at precisely $r$.
\end{proof}

\begin{theorem}[The whole original geometric strict ideal has the same strict drop]
\label{hir:actual-geometric-point-low-tangent-coefficient-strict-hs-drop}
\uses{hir:actual-point-low-tangent-coefficient-strict-hs-drop,hir:actual-original-point-normalized-geometric-strict-ideal}
\lean{Hironaka.HironakaActualPointTangentCoefficientHSDrop.hilbertSamuelValue_geometric_strict_lt_of_low_dehom_coefficient,Hironaka.HironakaActualPointTangentCoefficientHSDrop.dehom_coefficient_eq_zero_of_entire_point_hs_persistence}
\leanok
For an original analytic subset in dimension $n\ge2$, the same source
coefficient witness gives the strict finite drop for the entire actual
geometric strict-transform vanishing ideal in that point chart. In
particular, persistence of the entire original HS profile forces every
such pure dehomogenized coefficient below $\mu$ to vanish.
\end{theorem}
\begin{proof}
Apply the preceding result to the literal entire original origin ideal.
The genuine point-normalization/saturation/geometric-strict theorem
identifies the entire normalized ideal with the entire independently
defined geometric strict ideal, for this same actual map and carrier.
Substitution gives the strict inequality, without an incoming ideal
equality certificate. For the persistence assertion, a nonzero
coefficient would give this strict inequality at its finite degree,
contradicting the assumed equality there.
\end{proof}

Only the layer below the internally measured minimum degree is tested.
The all-vertex normalized-order criterion, current-history driver,
strict recursive invariant descent and termination remain separate.
